\section{Introduction}

% ¶1: Context — agents are powerful and widely deployed
AI coding agents have become powerful, widely deployed tools for software
development. They operate as long-running processes with access to shells,
source trees, package managers, and external APIs, autonomously writing code,
running tests, and managing infrastructure across extended
sessions~\cite{agentsight}.

% ¶2: Agents need harnesses to enforce behavioral contracts (consensus, state quickly)
As agents gain autonomy, projects govern them with behavioral contracts---do
not push directly to \texttt{main}; run tests before committing; never expose
secrets---encoded in instruction files such as \texttt{CLAUDE.md} and
\texttt{AGENTS.md}~\cite{chatlatanagulchai2025manifests,lulla2026agentsmd}.
Because LLM compliance is inherently probabilistic, deterministic enforcement
cannot live inside the model; it requires an external
harness~\cite{langchain-harness,fowler-harness}. The question is not whether
agents need harnesses, but where the harness should enforce and how it should
feed back.

% ¶3: Gap — single-level enforcement leaves an intent–behavior gap
Existing agent harness enforcement operates at a single level, leaving an
\emph{intent--behavior gap}. An agent's execution spans three abstraction
levels: \emph{intent} (what the agent declares it will and will not do),
\emph{action} (the tool calls it issues through its runtime), and
\emph{behavior} (the actual OS operations---\texttt{open}, \texttt{execve},
\texttt{connect}---that
result)~\cite{agentsight}. The mapping from action to behavior is many-to-many
and opaque: a single \texttt{run\_command("make")} can trigger hundreds of
system calls, and the same \texttt{git~push} can be reached through a direct
tool call, a \texttt{bash~-c} string, a Python subprocess, or a compiled
helper. Intent-level enforcement (prompt instructions) is probabilistic---the
agent may forget. Action-level enforcement (tool-call guards such as AgentSpec
and Progent~\cite{wang2025agentspec,shi2025progent}) is deterministic but
bypassable: it checks tool calls, not the OS behavior they produce.
Behavior-level enforcement (syscall filters, sandboxes) is unbypassable but
returns opaque \texttt{-EPERM} failures disconnected from intent. No single
level bridges the gap.

% ¶4: Insight — contracts are cross-object information flow
Real behavioral contracts frequently require cross-object information-flow
tracking that no single enforcement level provides. A contract like ``never
expose secrets to the network'' is not a predicate on one system call---it
requires knowing that the current process previously read a secret file and
now attempts a network connection. A contract like ``run tests before
committing'' requires knowing that a test process executed before a commit
process in the same session. A corpus study of instruction files from
50~popular projects confirms that such cross-object, cross-temporal contracts
are pervasive, yet today's agent
harnesses~\cite{langchain-harness,fowler-harness} keep their enforcement and
feedback components entirely at the application layer, where they cannot track
information flow across process, file, and network boundaries.

% ¶5: System — ActPlane bridges intent and behavior
ActPlane bridges intent and behavior by compiling agent-declared contracts into
OS-level labeled information-flow enforcement with corrective feedback. Agents
and developers declare contracts in a compact DSL---voluntary confinement,
analogous to \texttt{pledge()} in OpenBSD, where a program restricts its own
capabilities. ActPlane compiles these declarations into labeled
information-flow rules and loads them into an eBPF/BPF-LSM backend. Named
labels propagate across process, file, and network objects at kernel hooks,
encoding execution history as checkable state: when a process reads a file
labeled \texttt{SECRET}, the process acquires that label; when it forks, the
child inherits it. When a labeled subject violates a rule, ActPlane returns
corrective feedback to the agent---which contract was violated, why, and what
alternative path satisfies it---not \texttt{Permission~denied} but a reminder
of the agent's own declared intent.

% ¶6: Positioning — not new mechanism, unification at OS level
ActPlane is not a new enforcement mechanism but a unification at the OS level
of capabilities previously split across layers. CamQuery showed that
cross-channel label propagation and enforcement can be done in the
kernel~\cite{pasquier2018camquery}; AgentSpec demonstrated corrective feedback
for agents at the tool-call layer~\cite{wang2025agentspec}; FIDES and CaMeL
let agents participate in their own governance through application-layer
information-flow control~\cite{costa2025fides,debenedetti2025camel}. ActPlane
brings these together on the modern eBPF/BPF-LSM substrate, where contracts
survive context loss and cannot be bypassed below the tool API. It is not a
replacement for the application-layer harness but a complement---the
enforcement component that must reach below the tool API to the OS.

% ¶7: Results (TBD — fill with headline numbers when available)
% TODO: Add headline results here once evaluation is complete.
% Example: "ActPlane catches all five bypass paths that tool-layer guards miss,
% improves agent task-completion rate from X% to Y% under corrective feedback,
% incurs Z µs per-event overhead, and produces no false positives on N benign
% workloads."

% ¶8: Contributions
This paper makes three contributions.

\begin{enumerate}
\item It defines the \emph{intent--behavior gap} for AI agents and grounds it
with a corpus study of instruction files from 50~popular projects, showing
that cross-object information-flow contracts are pervasive in real agent
instruction files.

\item It presents ActPlane, a system that compiles agent-declared behavioral
contracts into eBPF/BPF-LSM labeled information-flow enforcement with
corrective feedback, propagating labels across process, file, and network
objects and reconnecting behavior-level violations to intent-level
remediation.

\item It evaluates ActPlane on bypass coverage across five tool paths, agent
recovery rate under four feedback conditions, false positives under benign
workloads, and per-event overhead, comparing against action-level and
behavior-level baselines.
\end{enumerate}
